{\itshape
The integration of wind power plants using doubly fed induction generator
(DFIG-WT) and permanent magnet synchronous generator (PMSG-WT) technologies
introduces distinct dynamic characteristics that may alter the small signal
stability (SSS) of power systems; this study analyzes and compares the SSS
impact of DFIG-WT and PMSG-WT integration on the IEEE nine-bus test system
through eigenvalue-based modal analysis using DIgSILENT PowerFactory 2024, with
six experimental scenarios designed by varying the wind turbine technology, the
integration bus location (Bus 7 and Bus 9), and the loading conditions. Results
show that DFIG-WT integration reduces the dominant mode damping ratio from 7.22\%
to approximately 4.04\% regardless of integration location, and significantly
narrows the load stability margin---the instability threshold for Load B drops
from 373\% to 86\%, and for Load A from 247\% to 63\%; by contrast, PMSG-WT
integration maintains or improves the damping ratio (7.25\%--8.73\%), as the
fully-rated converter decouples the machine dynamics from the network, keeping
the dominant oscillation mode governed by the remaining conventional synchronous
generators, and participation factor analysis confirms that the DFIG-WT shaft
variable ($\Delta\delta_{12}$) dominates the dominant oscillation mode while the
DFIG-WT internal rotor angle ($\varphi_m$) dominates the unstable mode under
load increase conditions. In conclusion, DFIG-WT and PMSG-WT exert opposing
effects on SSS---DFIG-WT degrades damping and narrows the safe operating margin,
while PMSG-WT maintains or improves stability with position-dependent
characteristics---providing technical guidance for planning wind power integration
that preserves the SSS of the power system.
}

\noindent\textbf{Keywords} : small signal stability, DFIG-WT, PMSG-WT, modal analysis, eigenvalue
